%=====================================================================
\chapter{Logistic Regression}\label{ch:logreg}
%=====================================================================
\locator{2}{Part II --- Linear Models and Generalization}

\begin{outcomes}
By the end of this chapter you will be able to:
\begin{tight}
  \item \textbf{Explain} why linear regression is unsuitable for classification.
  \item \textbf{Compute} a class probability with the sigmoid function, and
        \textbf{interpret} a weight as a change in log-odds and an odds ratio.
  \item \textbf{Find} and \textbf{draw} the decision boundary of a logistic
        regression model.
  \item \textbf{Compute} the cross-entropy loss and \textbf{perform} a gradient
        descent step by hand.
  \item \textbf{Extend} the model to several classes with one-vs-rest and
        softmax.
\end{tight}
\end{outcomes}

\begin{prereq}
\chref{linreg} (the linear model, gradient descent on a loss) and \chref{math}
(maximum likelihood, logarithms, the chain rule).
\end{prereq}

\begin{keyterms}
binary classification $\bullet$ sigmoid (logistic) function $\bullet$ odds
$\bullet$ log-odds (logit) $\bullet$ odds ratio $\bullet$ decision boundary
$\bullet$ threshold $\bullet$ cross-entropy (log loss) $\bullet$ linear
separability $\bullet$ one-vs-rest $\bullet$ softmax
\end{keyterms}

\section{Why Not Just Use Linear Regression?}

Code a binary label as $y = 1$ (late) or $y = 0$ (on time), fit a least-squares
line, and call anything above 0.5 late. It half works --- and it breaks in two
ways, both visible on the left of \dsref{linear-vs-logistic}. The line predicts
values below 0 and above 1, which cannot be probabilities. And a few examples far
from the boundary --- a small, well-specified project, which of course finished on
time --- drag the line towards them and move the 0.5 crossing, misclassifying
projects near the boundary that were previously correct. The model is being penalised for
being \emph{too} right about easy cases.

\dsfig{%
\begin{tikzpicture}
\begin{groupplot}[
  group style={group size=2 by 1, horizontal sep=1.4cm},
  width=6.8cm, height=4.8cm, axis lines=left,
  tick label style={font=\tiny}, label style={font=\tiny},
  title style={font=\scriptsize\bfseries},
  xlabel={risk score $x$}, ylabel={$P(\text{late})$ or prediction},
  xmin=0, xmax=10, ymin=-0.35, ymax=1.35]
\nextgroupplot[title={Least-squares line}]
\addplot[only marks, mark=*, mark size=1.5pt, greenhl!70!black] coordinates
  {(0.5,0)(1.2,0)(1.8,0)(2.5,0)(3.1,0)(3.8,0)(4.4,0)};
\addplot[only marks, mark=*, mark size=1.5pt, accent] coordinates
  {(5.2,1)(5.9,1)(6.6,1)(7.2,1)(7.9,1)(8.6,1)(9.4,1)};
\addplot[primary, line width=1.2pt, domain=0:10] {0.155*x-0.27};
\addplot[gray!50, dashed] coordinates {(0,0)(10,0)};
\addplot[gray!50, dashed] coordinates {(0,1)(10,1)};
\node[font=\tiny, accent, anchor=north west, align=left] at (axis cs:6.4,1.33)
     {above 1: not a\\probability};
\node[font=\tiny, accent, anchor=south west, align=left] at (axis cs:0.2,-0.33)
     {below 0};
\nextgroupplot[title={Logistic (sigmoid) curve}]
\addplot[only marks, mark=*, mark size=1.5pt, greenhl!70!black] coordinates
  {(0.5,0)(1.2,0)(1.8,0)(2.5,0)(3.1,0)(3.8,0)(4.4,0)};
\addplot[only marks, mark=*, mark size=1.5pt, accent] coordinates
  {(5.2,1)(5.9,1)(6.6,1)(7.2,1)(7.9,1)(8.6,1)(9.4,1)};
\addplot[primary, line width=1.2pt, domain=0:10, samples=100] {1/(1+exp(-2.2*(x-4.8)))};
\addplot[gray!50, dashed] coordinates {(0,0)(10,0)};
\addplot[gray!50, dashed] coordinates {(0,1)(10,1)};
\addplot[gray!50, dashed] coordinates {(4.8,-0.35)(4.8,0.5)};
\node[font=\tiny, primary, anchor=west] at (axis cs:5.0,0.3) {$P = 0.5$ here};
\end{groupplot}
\end{tikzpicture}}%
{The same on-time (green) and late (red) projects modelled two ways. The straight line
leaves $[0,1]$ and is pulled by easy cases; the sigmoid stays between 0 and 1 and
flattens out where the answer is already clear.}{linear-vs-logistic}

\section{The Sigmoid and the Log-Odds}

\begin{definitionbox}[Logistic Regression]
Compute a linear score $z = \mathbf{w}\cdot\mathbf{x} + b$ exactly as in linear
regression, then squash it into $(0, 1)$ with the \term{sigmoid} (logistic)
function:
\[
P(y = 1 \mid \mathbf{x}) = \sigma(z) = \frac{1}{1 + e^{-z}}
\]
$\sigma(0) = 0.5$; large positive $z$ gives a probability near 1 and large negative
$z$ a probability near 0. Predict class 1 when the probability exceeds a
\term{threshold}, by default 0.5.
\end{definitionbox}

Inverting the sigmoid shows what the linear score means. If $p = \sigma(z)$ then
\[
z = \ln\frac{p}{1-p}
\]
The ratio $p/(1-p)$ is the \term{odds} --- a probability of 0.8 is odds of 4 to 1
--- and its logarithm is the \term{log-odds} or \term{logit}. So logistic
regression is a linear model \emph{of the log-odds}. Each weight is the change in
log-odds for a one-unit increase in its feature, and $e^{w_j}$ is the
\term{odds ratio}: the factor by which the odds are multiplied.

\begin{worked}[Probabilities and Odds for Two Projects]
A fitted model of finishing late: $z = 4 - 5\,c - 3\,s$, where $c$ is requirement
clarity (the fraction of requirements signed off) and $s$ the skill match (the
fraction of the needed skills the team already has), both on a 0--1 scale.

\smallskip
\textbf{Project A} ($c = 0.9$, $s = 0.8$):
$z = 4 - 4.5 - 2.4 = -2.9$, so
$P(\text{late}) = \dfrac{1}{1 + e^{2.9}} = \dfrac{1}{1 + 18.17} = \mathbf{0.052}$.

\textbf{Project B} ($c = 0.5$, $s = 0.4$):
$z = 4 - 2.5 - 1.2 = 0.3$, so
$P(\text{late}) = \dfrac{1}{1 + e^{-0.3}} = \dfrac{1}{1 + 0.741} = \mathbf{0.574}$.

\smallskip
\textbf{Reading the clarity weight.} $w_c = -5$ per unit of clarity, so ten
percentage points more clarity ($+0.1$) changes the log-odds by $-0.5$ and
multiplies the odds of finishing late by $e^{-0.5} = \mathbf{0.61}$ --- a 39\%
reduction in the odds, whatever the project's starting point. Note: the
\emph{odds}, not the probability. For Project A the probability barely moves; for
Project B, near the middle of the curve, it moves a lot.
\end{worked}

\section{The Decision Boundary}

The model predicts class 1 when $\sigma(z) > 0.5$, which happens exactly when
$z > 0$. So the \term{decision boundary} is the set of points where
$\mathbf{w}\cdot\mathbf{x} + b = 0$: a straight line in two dimensions, a plane in
three, a hyperplane in general. Logistic regression is a \emph{linear
classifier}; it can separate classes only with a flat boundary.

\dsfig{%
\begin{tikzpicture}
\begin{axis}[width=8.2cm, height=6.4cm, axis lines=left,
  xlabel={requirement clarity $c$}, ylabel={skill match $s$},
  tick label style={font=\tiny}, label style={font=\scriptsize},
  xmin=0.2, xmax=1.0, ymin=0, ymax=1.0]
\addplot[name path=b, primary, line width=1.3pt, domain=0.2:1.0] {(4-5*x)/3};
\addplot[primary!60, dashed, line width=0.8pt, domain=0.2:1.0] {(4-1.0986-5*x)/3};
\addplot[primary!60, dashed, line width=0.8pt, domain=0.2:1.0] {(4+1.0986-5*x)/3};
\addplot[only marks, mark=*, mark size=1.5pt, accent] coordinates
  {(0.30,0.35)(0.35,0.70)(0.42,0.20)(0.46,0.52)(0.55,0.30)(0.38,0.48)(0.60,0.12)(0.50,0.40)};
\addplot[only marks, mark=*, mark size=1.5pt, greenhl!70!black] coordinates
  {(0.70,0.55)(0.85,0.40)(0.92,0.75)(0.78,0.85)(0.66,0.95)(0.95,0.20)(0.88,0.62)(0.74,0.30)};
\node[font=\tiny, primary, fill=white, inner sep=1pt, rotate=-44] at (axis cs:0.62,0.30)
     {$z = 0$: $P(\text{late}) = 0.5$};
\node[font=\tiny, primary!70, anchor=west, fill=white, inner sep=1pt] at (axis cs:0.43,0.96)
     {$P(\text{late}) = 0.25$};
\node[font=\tiny, primary!70, anchor=west, fill=white, inner sep=1pt] at (axis cs:0.21,0.66)
     {$P(\text{late}) = 0.75$};
\node[font=\tiny\bfseries, accent, anchor=west] at (axis cs:0.22,0.08) {predicted late};
\node[font=\tiny\bfseries, greenhl!55!black, anchor=east] at (axis cs:0.99,0.95) {predicted on time};
\end{axis}
\end{tikzpicture}}%
{The decision boundary of the model $z = 4 - 5c - 3s$ is the line $5c + 3s = 4$. The
dashed lines, where $z = \pm\ln 3$, are where the probability of finishing late is 0.75
and 0.25: probability changes gradually across the boundary.}{logreg-boundary}

\section{The Loss: Cross-Entropy}

Squared error is the wrong loss for probabilities: combined with the sigmoid it
gives a bumpy surface with flat regions where gradient descent stalls. The right
loss comes, again, from maximum likelihood. If the model gives probability $p_i$ to
class 1, the likelihood of the observed label is $p_i$ when $y_i = 1$ and
$1 - p_i$ when $y_i = 0$. Taking minus the average log of that gives:

\begin{definitionbox}[Cross-Entropy (Log Loss)]
\[
L(\mathbf{w}, b) = -\frac{1}{n}\sum_{i=1}^{n}\Big[\,y_i \ln p_i + (1 - y_i)\ln(1 - p_i)\,\Big],
\qquad p_i = \sigma(\mathbf{w}\cdot\mathbf{x}_i + b)
\]
For each example only one term survives: $-\ln p_i$ if the label is 1,
$-\ln(1 - p_i)$ if it is 0. The loss is convex --- one bowl, one minimum ---
so gradient descent finds the best weights from any start.
\end{definitionbox}

\dsfig{%
\begin{tikzpicture}
\begin{axis}[width=8.0cm, height=4.8cm, axis lines=left,
  xlabel={predicted probability of class 1, $p$}, ylabel={loss for one example},
  tick label style={font=\tiny}, label style={font=\scriptsize},
  xmin=0, xmax=1, ymin=0, ymax=4.8,
  legend style={font=\tiny, draw=gray!40, at={(0.5,0.98)}, anchor=north}]
\addplot[greenhl!70!black, line width=1.3pt, domain=0.009:1, samples=120] {-ln(x)};
\addlegendentry{true label 1: $-\ln p$}
\addplot[accent, line width=1.3pt, domain=0:0.991, samples=120] {-ln(1-x)};
\addlegendentry{true label 0: $-\ln(1-p)$}
\node[font=\tiny, greenhl!45!black, anchor=west, align=left] at (axis cs:0.07,2.9)
     {confidently wrong:\\$p = 0.01$ costs 4.6};
\draw[gray!55, line width=0.4pt] (axis cs:0.5,0.74) -- (axis cs:0.5,1.75);
\node[font=\tiny, gray!55!black, anchor=south] at (axis cs:0.5,1.75) {unsure: 0.69 either way};
\end{axis}
\end{tikzpicture}}%
{Cross-entropy punishes confident mistakes without limit. Being right at 0.9 costs
0.11; being wrong at 0.99 costs 4.6.}{log-loss}

\begin{worked}[Cross-Entropy on Four Predictions]
\begin{tight}
  \item Label 1, predicted $p = 0.9$: loss $-\ln 0.9 = 0.105$.
  \item Label 1, predicted $p = 0.6$: loss $-\ln 0.6 = 0.511$.
  \item Label 0, predicted $p = 0.2$: loss $-\ln 0.8 = 0.223$.
  \item Label 0, predicted $p = 0.7$: loss $-\ln 0.3 = 1.204$.
\end{tight}
\textbf{Average:} $(0.105 + 0.511 + 0.223 + 1.204)/4 = \mathbf{0.511}$.

\smallskip
The last example is the only one misclassified at a 0.5 threshold, and it
contributes more than half the total. Replace it by a label-1 example predicted at
$p = 0.01$ and that one example alone costs $-\ln 0.01 = 4.6$. Cross-entropy cares
not just whether you are wrong, but how sure you were.
\end{worked}

\section{Fitting by Gradient Descent}

Differentiating the cross-entropy through the sigmoid, most of the algebra
cancels and leaves a result of exactly the same shape as linear regression's:
\[
\frac{\partial L}{\partial w_j} = \frac{1}{n}\sum_{i=1}^{n}(p_i - y_i)\,x_{ij},
\qquad
\frac{\partial L}{\partial b} = \frac{1}{n}\sum_{i=1}^{n}(p_i - y_i)
\]
Error times feature, averaged --- with the error now measured in probability.

\begin{worked}[One Gradient Descent Step]
Four examples with one feature: $x = 1, 2, 3, 4$ and labels $y = 0, 0, 1, 1$. Start at
$w = 0$, $b = 0$ with learning rate $\alpha = 1$.

\smallskip
\textbf{Predictions:} $z = 0$ everywhere, so $p = 0.5$ for all four, and the loss is
$-\ln 0.5 = 0.693$.

\textbf{Errors} $p - y$: $0.5,\ 0.5,\ -0.5,\ -0.5$.
\[
\frac{\partial L}{\partial w} = \tfrac{1}{4}\big[(0.5)(1) + (0.5)(2) + (-0.5)(3) + (-0.5)(4)\big]
= \tfrac{1}{4}(-2) = -0.5,
\qquad \frac{\partial L}{\partial b} = 0
\]
\textbf{Update:} $w \leftarrow 0 - 1 \times (-0.5) = 0.5$; $b$ is unchanged.

\textbf{New probabilities} $\sigma(0.5x)$: $0.62,\ 0.73,\ 0.82,\ 0.88$. The loss falls to
$0.654$. The two class-1 examples are already right; the next steps will push $b$
negative so that the first two fall below 0.5 as well.
\end{worked}

\begin{alertbox}[When the Classes Can Be Separated Perfectly]
In the worked example a single threshold, $x = 2.5$, separates the classes with no
errors. The loss can then always be reduced by making the sigmoid steeper, so the
weights grow without limit: after 5{,}000 steps $w \approx 8.9$ and still rising,
while the boundary settles at $x \approx 2.5$. The fix is regularisation
(\chref{generalization}), which scikit-learn applies by default --- that is what
its parameter \texttt{C} controls.
\end{alertbox}

\section{More Than Two Classes}

\textbf{One-vs-rest} trains one binary classifier per class (``class $k$ or not?'')
and predicts the class whose classifier is most confident. \textbf{Softmax}
(multinomial) regression is the direct generalisation: one linear score per class,
turned into probabilities that sum to 1:
\[
P(y = k \mid \mathbf{x}) = \frac{e^{z_k}}{\sum_{j} e^{z_j}},
\qquad z_k = \mathbf{w}_k\cdot\mathbf{x} + b_k
\]

\begin{worked}[Softmax on Three Scores]
A ticket classifier gives scores $z = (2.0,\ 1.0,\ 0.1)$ for the classes
\emph{bug}, \emph{feature request} and \emph{question}.
\[
e^{2.0} = 7.389,\quad e^{1.0} = 2.718,\quad e^{0.1} = 1.105,\qquad \text{sum} = 11.212
\]
\[
P = \left(\frac{7.389}{11.212},\ \frac{2.718}{11.212},\ \frac{1.105}{11.212}\right)
  = (\mathbf{0.659},\ \mathbf{0.242},\ \mathbf{0.099})
\]
Adding the same constant to every score leaves the probabilities unchanged ---
only the differences between scores matter. The same function is the output layer
of almost every neural network classifier (\chref{ann}).
\end{worked}

\begin{pitfall}
\begin{tight}
  \item \textbf{Reading a weight as a change in probability.} It is a change in
        log-odds. The effect on probability depends on where on the curve you
        start.
  \item \textbf{Assuming the 0.5 threshold is right.} It is a default. When the
        two errors cost different amounts, move it (\chref{evaluation}).
  \item \textbf{Forgetting that the boundary is linear.} If the classes are
        separated by a curve, add features (squares, products) or use a
        non-linear model (Part III).
  \item \textbf{Comparing weights of unscaled features.} As in \chref{linreg},
        standardise first.
  \item \textbf{Trusting probabilities from an unregularised model on separable
        data.} They will all be pushed to 0 or 1.
\end{tight}
\end{pitfall}

%---------------------------------------------------------------------
\begin{lab}[Logistic Regression by Hand and by Library]
\begin{lstlisting}[language=Python]
import numpy as np
from sklearn.model_selection import train_test_split
from sklearn.pipeline import make_pipeline
from sklearn.preprocessing import StandardScaler
from sklearn.linear_model import LogisticRegression

sigmoid = lambda z: 1 / (1 + np.exp(-z))

# --- 1. the worked example: gradient descent by hand ---------------------
x = np.array([1, 2, 3, 4.0])
y = np.array([0, 0, 1, 1.0])
logloss = lambda p: -np.mean(y * np.log(p) + (1 - y) * np.log(1 - p))
w, b = 0.0, 0.0
p = sigmoid(w * x + b)
print("before: loss", round(logloss(p), 3))
w -= 1.0 * np.mean((p - y) * x)          # gradient: mean of (p - y) * x
b -= 1.0 * np.mean(p - y)
print(f"after one step: w = {w}, b = {b}, "
      f"loss {logloss(sigmoid(w * x + b)):.3f}")
for step in range(5000):                 # ... and many more steps
    p = sigmoid(w * x + b)
    w, b = w - np.mean((p - y) * x), b - np.mean(p - y)
print(f"after 5000 steps: w = {w:.2f}, b = {b:.2f}, "
      f"boundary at x = {-b / w:.2f}")

# --- 2. a real classifier, read as odds ratios ---------------------------
rng = np.random.default_rng(3)
n = 500                                  # past projects
names = ["team_size", "changes", "dependencies", "clarity"]
X = np.column_stack([rng.integers(3, 16, n),          # people
                     rng.poisson(5, n),                # change requests
                     rng.integers(0, 7, n),            # external teams
                     rng.uniform(0.3, 1.0, n)])        # clarity, 0-1
z = -1.5 + 0.05 * X[:, 0] + 0.3 * X[:, 1] + 0.4 * X[:, 2] - 3.5 * X[:, 3]
y = (rng.random(n) < 1 / (1 + np.exp(-z))).astype(int)   # 1 = late
X_tr, X_te, y_tr, y_te = train_test_split(
    X, y, test_size=0.25, random_state=0, stratify=y)
clf = make_pipeline(StandardScaler(), LogisticRegression())
clf.fit(X_tr, y_tr)
print("test accuracy:", round(clf.score(X_te, y_te), 3))
coef = clf[-1].coef_[0]
for i in np.argsort(-np.abs(coef)):      # strongest risk factors first
    print(f"  {names[i]:13s} w = {coef[i]:5.2f}  "
          f"odds x {np.exp(coef[i]):.2f} per standard deviation")
print("P(late) for 3 test projects:",
      clf.predict_proba(X_te[:3])[:, 1].round(3))

# --- 3. three classes: softmax -------------------------------------------
m = 300                                  # support tickets
T = np.column_stack([rng.uniform(0, 4, m),            # log10 users hit
                     rng.integers(0, 2, m)])           # 1 = workaround
score = T[:, 0] - 1.2 * T[:, 1] + rng.normal(0, 0.4, m)
level = np.digitize(score, [1.0, 2.5])   # 0 low, 1 medium, 2 high
soft = make_pipeline(StandardScaler(), LogisticRegression())
soft.fit(T, level)
print("P(low, medium, high) for ticket 60:",
      soft.predict_proba(T[[60]]).round(3))
\end{lstlisting}
In part 1, change the range to 50{,}000 steps. Does the boundary move? Do the
weights? Explain both answers.
\end{lab}

%---------------------------------------------------------------------
\begin{checkpoint}
\begin{tightnum}
  \item A model gives $z = 1.2$ for a task. What is the probability that it
        finishes late, and what are the odds?
  \item A weight is $+0.7$. By what factor does a one-unit increase in its feature
        multiply the odds?
  \item Which costs more under cross-entropy: predicting 0.4 for a positive
        example, or 0.4 for a negative one? Compute both.
\end{tightnum}
\end{checkpoint}

%---------------------------------------------------------------------
\begin{chaptersummary}
\begin{tight}
  \item Linear regression is unsuitable for classification: its outputs leave
        $[0,1]$ and easy cases distort its boundary.
  \item Logistic regression passes a linear score through the sigmoid,
        $p = 1/(1 + e^{-z})$; the score is the log-odds of class 1.
  \item Each weight changes the log-odds by $w$ per unit; $e^w$ is the odds ratio.
  \item The decision boundary $\mathbf{w}\cdot\mathbf{x} + b = 0$ is linear.
  \item Cross-entropy is the maximum likelihood loss; it is convex and punishes
        confident mistakes heavily. Its gradient is the mean of
        $(p - y) \times$ feature.
  \item Softmax extends the model to $k$ classes; on separable data the weights
        grow without limit unless regularised.
\end{tight}
\end{chaptersummary}

\begin{reviewq}
  \item \textbf{(MCQ)} $\sigma(0)$ equals: (a)~0 (b)~0.5 (c)~1 (d)~$e$.
  \item \textbf{(MCQ)} The decision boundary of logistic regression is:
        (a)~a sigmoid curve (b)~a hyperplane (c)~a circle (d)~any shape.
  \item \textbf{(MCQ)} A probability of 0.75 corresponds to odds of: (a)~0.75
        (b)~1.33 (c)~3 (d)~4.
  \item \textbf{(Short)} Explain why cross-entropy rather than squared error is
        used to train logistic regression.
  \item \textbf{(Short)} What happens to logistic regression's weights on data
        that a straight line separates perfectly, and why?
  \item \textbf{(Short)} Contrast one-vs-rest with softmax regression.
  \item \textbf{(Applied)} With $z = -3 + 2x_1 + 0.5x_2$, compute the probability
        of class 1 for $(x_1, x_2) = (1, 2)$ and $(2, 0)$. Write the equation of
        the decision boundary and state which side each point lies on.
  \item \textbf{(Applied)} Perform a second gradient descent step on the worked
        example, starting from $w = 0.5$, $b = 0$, with $\alpha = 1$.
\end{reviewq}
